\section{Colophon: how this document is built}
\label{sec:colophon}

The toolchain is recorded because the paper claims reproducibility, and a
reproducibility claim that omits the build is incomplete.

\subsection*{Files}

\begin{center}\small
\begin{tabularx}{\textwidth}{@{}p{5.0cm} X@{}}
\toprule
\file{ontodag-fca.tex} & the document: preamble, styles, and the
  \verb|\input| list\\
\file{sections/*.tex} & one file per section, in reading order\\
\file{ontodag-fca.bib} & the bibliography, grouped by literature\\
\file{experiments/fca_experiments.py} & every number in the paper\\
\file{experiments/results.txt} & its output (Appendix~\ref{sec:experiments})\\
\file{experiments/explosion.dat} & the data behind
  Figure~\ref{fig:explosion}\\
\file{figures/*.dot} & Graphviz sources; four generated by the script,
  three hand-written\\
\file{figures/*.tex} & TikZ fragments produced from them by
  \texttt{dot2tex}\\
\texttt{Makefile} & the build\\
\bottomrule
\end{tabularx}
\end{center}

\subsection*{The build}

\begin{lstlisting}
$ make experiments   # rerun the measurements and regenerate the .dot files
$ make figures       # dot2tex: figures/*.dot -> figures/*.tex
$ make               # pdflatex, bibtex, pdflatex, pdflatex
\end{lstlisting}

Requirements: \texttt{pdflatex} and \texttt{bibtex} (TeX~Live 2023 or later,
with \texttt{tikz}, \texttt{pgfplots}, \texttt{tcolorbox}, \texttt{natbib},
\texttt{stmaryrd}), \texttt{graphviz} for \texttt{dot}, \texttt{dot2tex}, and
Python 3 with \odag{} importable (the script inserts \texttt{../src} on the
path, so a checkout suffices; no installation needed). The certificate
experiment additionally needs \texttt{recordstore}, and skips with a printed
note if it is absent.

\subsection*{How the figures are made}

Two kinds, deliberately.

\paragraph{Graph figures come from Graphviz.} A \texttt{.dot} file is laid out
by \texttt{dot} and converted to a TikZ fragment by \texttt{dot2tex}
(\verb|--format tikz --codeonly --texmode raw|), which emits bare
\verb|\node| and \verb|\draw| commands with coordinates. The paper wraps each
fragment in its own \texttt{tikzpicture} via the \verb|\dotfig| macro, so
scale, fonts and node styles are the document's decision rather than
Graphviz's. Two consequences worth recording for anyone reusing the setup:

\begin{itemize}[leftmargin=1.4em]
\item In \texttt{raw} texmode a label is LaTeX, so \verb|&| must be avoided or
  escaped and a line break must reach LaTeX as \verb|\\| --- which means four
  backslashes in the \texttt{.dot} source. The wrapper sets
  \texttt{align=center} on every node so that \verb|\\| is legal.
\item Graphviz colour attributes pass through as raw names, and HSB triples
  (Graphviz's default numeric colour syntax) are not a colour model
  \textsc{pgf} understands. The figures therefore use no colour at all and
  distinguish node roles by shape and border count --- which also means they
  survive greyscale printing.
\end{itemize}

\paragraph{Conceptual figures are hand-written TikZ}, because their content is
an argument rather than data: Figures~\ref{fig:roadmap},
\ref{fig:hasse}, \ref{fig:mind-brain}, \ref{fig:verify} and
\ref{fig:spectrum}. The chart in Figure~\ref{fig:explosion} is
\texttt{pgfplots} reading the generated \texttt{.dat} file, so it cannot
disagree with the text.

\subsection*{Relationship to the companion paper}

A companion manuscript in the same directory, \texttt{main.tex}
\citep{ontodagswarm}, presents \odag{} as a semantic layer for distributed
storage: the data structure, the persistence mapping onto Swarm, the
multi-writer merge, and the MDL learning direction. It overlaps this one in
Sections~\ref{sec:ontodag} and~\ref{sec:crypto} and is the better starting
point for a reader who wants the systems story rather than the order theory.
This paper is the theoretical companion: what the structure \emph{is}, in the
terms of the field that has studied such structures for forty years.

\subsection*{Provenance of this document}

Drafted in August 2026 with Claude (Anthropic) as a writing and analysis
collaborator, working from the \odag{} repository's design records --- in
particular the notes on semantic codes, the higher-layer contract, the
database-direction walls, and the historical survey of philosophical
languages, all of which are cited in the body where their arguments are used.
The experiments were written for this paper and run before the claims that
quote them; the transcripts in Sections~\ref{sec:ontodag}
and~\ref{sec:agents} are real sessions, not illustrations.
